% !TeX program = xelatex
\documentclass[12pt,a4paper]{article}
\usepackage[margin=2cm]{geometry}
\usepackage{fontspec}
\usepackage{amsmath,amssymb}
\usepackage{enumitem}
\setmainfont{Times New Roman}
\setlength{\parindent}{0pt}
\setlength{\parskip}{6pt}
\begin{document}
\begin{center}
{\Large\bfseries ĐÁP ÁN ĐỀ THI TUYỂN SINH VÀO LỚP 10 THPT}\\
{\bfseries Thành phố Hà Nội, năm học 2011--2012}\\
{\itshape Phần Đại số}
\end{center}

\textbf{Câu 1.}
\begin{enumerate}[label=\alph*)]
\item Với $x\ge0$, $x\ne25$,
\begin{align*}
A&=\frac{\sqrt{x}(\sqrt{x}+5)-10\sqrt{x}-5(\sqrt{x}-5)}
{(\sqrt{x}-5)(\sqrt{x}+5)}\\
&=\frac{x-10\sqrt{x}+25}{x-25}
=\frac{(\sqrt{x}-5)^2}{(\sqrt{x}-5)(\sqrt{x}+5)}
=\frac{\sqrt{x}-5}{\sqrt{x}+5}.
\end{align*}

\item Với $x=9$,
\[
A=\frac{3-5}{3+5}=-\frac14.
\]

\item Vì $\sqrt{x}+5>0$,
\begin{align*}
\frac{\sqrt{x}-5}{\sqrt{x}+5}<\frac13
&\Longleftrightarrow 3\sqrt{x}-15<\sqrt{x}+5\\
&\Longleftrightarrow \sqrt{x}<10\\
&\Longleftrightarrow 0\le x<100.
\end{align*}
Kết hợp điều kiện xác định:
\[
\boxed{0\le x<100,\quad x\ne25}.
\]
\end{enumerate}

\textbf{Câu 2.}

Gọi $a$ là số tấn hàng dự định chở mỗi ngày và $b$ là số ngày theo kế hoạch ($a>0$, $b\in\mathbb N^*$). Ta có
\[
\begin{cases}
ab=140,\\
(a+5)(b-1)=150.
\end{cases}
\]
Từ phương trình thứ hai và $ab=140$:
\[
5b-a=15.
\]
Thay $a=\dfrac{140}{b}$:
\[
5b-\frac{140}{b}=15
\Longleftrightarrow b^2-3b-28=0
\Longleftrightarrow (b-7)(b+4)=0.
\]
Do $b>0$, suy ra $b=7$. Vậy theo kế hoạch đội xe chở hàng trong
$\boxed{7\text{ ngày}}$.

\textbf{Câu 3.}
\begin{enumerate}[label=\alph*)]
\item Khi $m=1$, phương trình hoành độ giao điểm là
\[
x^2=2x+8
\Longleftrightarrow (x+2)(x-4)=0.
\]
Với $x=-2$ thì $y=4$; với $x=4$ thì $y=16$. Vậy hai giao điểm là
\[
\boxed{A(-2;4),\quad B(4;16)}.
\]

\item Phương trình hoành độ giao điểm là
\[
x^2-2x+m^2-9=0. \tag{1}
\]
Hai giao điểm nằm về hai phía của trục tung khi và chỉ khi (1) có hai nghiệm trái dấu:
\[
x_1x_2=m^2-9<0
\Longleftrightarrow -3<m<3.
\]
Vậy $\boxed{-3<m<3}$.
\end{enumerate}

\textbf{Câu 5.}

Với $x>0$,
\begin{align*}
M
&=4\left(x-\frac12\right)^2+x+\frac1{4x}+2010\\
&\ge 0+2\sqrt{x\cdot\frac1{4x}}+2010=2011.
\end{align*}
Dấu bằng xảy ra khi $x=\dfrac12$. Vậy
\[
\boxed{\min M=2011}.
\]
\end{document}
